\Needspace{6\baselineskip}
\begin{center}\hypertarget{native-252-article180}{}\pdfbookmark[2]{180.}{article-180}180.\end{center}
Principium in art. praec. expositum eo quoque nomine se commendat, quod
determinatio incognitarum numerica ad algorithmum expeditissimum reducitur,
quoties functiones $V$, $V'$, $V''$ etc. lineares sunt. Supponamus esse
\[
\begin{gathered}
M-V=v=-m+ap+bq+cr+ds+\text{etc.}\\
M'-V'=v'=-m'+a'p+b'q+c'r+d's+\text{etc.}\\
M''-V''=v''=-m''+a''p+b''q+c''r+d''s+\text{etc.}\\
\text{etc.,}
\end{gathered}
\]
statuamusque
\[
\begin{gathered}
av+a'v'+a''v''+\text{etc.}=P\\
bv+b'v'+b''v''+\text{etc.}=Q\\
cv+c'v'+c''v''+\text{etc.}=R\\
dv+d'v'+d''v''+\text{etc.}=S\\
\text{etc.}
\end{gathered}
\]
Tunc $\nu$ aequationes art. 177, e quibus incognitarum valores determinare
oportet, manifesto hae erunt:
\[
P=0,\qquad Q=0,\qquad R=0,\qquad S=0\quad\text{etc.,}
\]
siquidem observationes aeque bonas supponimus, ad quem casum reliquos reducere
in art. praec. docuimus. Adsunt itaque totidem aequationes lineares, quot
incognitae determinandae sunt, unde harum valores per eliminationem vulgarem
elicientur.

Videamus nunc, utrum haec eliminatio semper possibilis sit, an unquam
\SourceLeaf{253}solutio indeterminata vel adeo impossibilis evadere possit. Ex
eliminationis theoria constat, casum secundum vel tertium tunc locum habiturum
esse, quando ex aequationibus $P=0$, $Q=0$, $R=0$, $S=0$ etc., omissa una,
aequatio conflari potest vel identica cum omissa vel eidem repugnans, sive,
quod eodem redit, quando assignare licet functionem linearem
\[
\alpha P+\beta Q+\gamma R+\delta S+\text{etc.,}
\]
quae fit identice vel $=0$ vel saltem ab omnibus incognitis $p$, $q$, $r$, $s$
etc. libera. Supponamus itaque fieri
$\alpha P+\beta Q+\gamma R+\delta S+\text{etc.}=\varkappa$. Sponte habetur
aequatio identica
\[
\begin{aligned}
&(v+m)v+(v'+m')v'+(v''+m'')v''+\text{etc.}\\
&\qquad=pP+qQ+rR+sS+\text{etc.}
\end{aligned}
\]
Quodsi itaque per substitutiones $p=\alpha x$, $q=\beta x$, $r=\gamma x$,
$s=\delta x$ etc. functiones $v$, $v'$, $v''$ etc. resp. in $-m+\lambda x$,
$-m'+\lambda'x$, $-m''+\lambda''x$ etc. transire supponimus, manifesto aderit
aequatio identica
\[
\begin{aligned}
&(\lambda\lambda+\lambda'\lambda'+\lambda''\lambda''+\text{etc.})xx\\
&\qquad-(\lambda m+\lambda'm'+\lambda''m''+\text{etc.})x=\varkappa x,
\end{aligned}
\]
i. e. erit $\lambda\lambda+\lambda'\lambda'+\lambda''\lambda''+\text{etc.}=0$,
$\varkappa+\lambda m+\lambda'm'+\lambda''m''+\text{etc.}=0$: hinc vero
necessario esse debebit $\lambda=0$, $\lambda'=0$, $\lambda''=0$ etc. atque
$\varkappa=0$. Hinc patet, functiones omnes $V$, $V'$, $V''$ etc. ita
comparatas esse, ut valores ipsarum non mutentur, si quantitates $p$, $q$, $r$,
$s$ etc. capiant incrementa vel decrementa quaecunque numeris $\alpha$,
$\beta$, $\gamma$, $\delta$ etc. proportionalia: huiusmodi autem casus, in
quibus manifesto determinatio incognitarum ne tunc quidem possibilis esset,
si ipsi veri valores functionum $V$, $V'$, $V''$ etc. darentur, huc non
pertinere iam supra monuimus.

Ceterum ad casum hic consideratum omnes reliquos, ubi functiones $V$, $V'$,
$V''$ etc. non sunt lineares, facile reducere possumus. Scilicet designantibus
$\pi$, $\chi$, $\rho$, $\sigma$ etc. valores approximatos incognitarum $p$,
$q$, $r$, $s$ etc. (quos facile eliciemus, si ex $\mu$ aequationibus $V=M$,
$V'=M'$, $V''=M''$ etc. primo $\nu$ tantum in usum vocamus), introducemus
incognitarum loco alias $p'$, $q'$, $r'$, $s'$ etc., statuendo $p=\pi+p'$,
$q=\chi+q'$, $r=\rho+r'$, $s=\sigma+s'$ etc.: manifesto harum novarum
incognitarum valores tam parvi erunt, ut quadrata productaque negligere
liceat, quo pacto aequationes sponte evadent lineares. Quodsi dein calculo
absoluto contra exspectationem valores incognitarum $p'$, $q'$, $r'$, $s'$
etc. \SourceLeaf{254}tanti emergerent, ut parum tutum videatur, quadrata
productaque neglexisse, eiusdem operationis repetitio (acceptis loco ipsarum
$\pi$, $\chi$, $\rho$, $\sigma$ etc. valoribus correctis ipsarum $p$, $q$,
$r$, $s$ etc.) remedium promtum afferet.
